%---
% title: "PowerShell Cheatsheet"
% date: 2026-08-18
% tags: [powershell, windows]
% summary: "Tổng hợp nhanh PowerShell từ cơ bản đến pipeline, provider, script và error handling."
%---
\section{1. Getting Started}
\textbf{PowerShell} = a command line shell + a scripting language. On Windows 7+, and on Linux/macOS via \emph{PowerShell Core}.
\begin{itemize}
\item Open it: Start menu, type \texttt{PowerShell}.
\item \textbf{Run as Administrator}: right-click, Run as Administrator (needed for the registry, system files, \texttt{Update-Help}).
\item The \texttt{PS>} prompt means PowerShell is ready for input.
\item \textbf{TAB} cycles completions forward, \textbf{SHIFT-TAB} backward. Works for command names and parameters.
\end{itemize}

\textbf{Familiar DOS commands}

\begin{verbatim}
cd .\Windows\  # change directory
dir  # list contents
cls  # clear the screen
Get-Alias  # show every alias
\end{verbatim}
Those DOS commands are really \textbf{aliases} pointing at PowerShell cmdlets.

\section{2. Commands}
Command names follow the \textbf{Verb-Noun} convention: \texttt{Get-Content}, \texttt{Set-Variable}.
\begin{itemize}
\item \textbf{Cmdlet}: written in C\#, shipped by Microsoft.
\item \textbf{Function}: written in PowerShell (these are what you will write).
\item Also \textbf{aliases} and \textbf{external scripts}.
\end{itemize}

\textbf{Get-Command}

\begin{verbatim}
Get-Command  # everything
Get-Command -Verb Get
Get-Command -Noun Content
Get-Command -Verb Get -Noun Content
Get-Command -Name Get-Content
\end{verbatim}
Common verbs: \texttt{Get} (read), \texttt{Set}, \texttt{Update}, \texttt{Remove} (modify).

\section{3. Help}
\begin{verbatim}
Get-Help Add-Content
Get-Help Add-Content -Examples  # most useful
Get-Help Add-Content -Detailed
Get-Help Add-Content -Full
Get-Help about_Core_Commands
Get-Help -Name About*  # wildcard *
Update-Help  # needs Admin
\end{verbatim}
\begin{itemize}
\item Help sections: SYNOPSIS, SYNTAX, DESCRIPTION, RELATED LINKS, REMARKS.
\item \textbf{About topics}: broad subjects, not tied to a single command.
\item \textbf{Positional parameter}: drop the parameter name and pass the value by position.
\end{itemize}

\section{4. Variables}
Every variable starts with \texttt{\$}.
\begin{verbatim}
$color = 'blue'  # declare + assign
$color  # read the value
Set-Variable -Name color -Value blue
Get-Variable -Name color
Get-Variable  # all variables
\end{verbatim}
Referencing an undefined variable raises \emph{VariableIsUndefined}.

\textbf{Strict mode}

\begin{verbatim}
Set-StrictMode -Version Latest
\end{verbatim}
Throws on undefined variables and properties. Best practice when writing scripts.

\textbf{Automatic variables}

Created by PowerShell -- treat them as read-only:
\begin{itemize}
\item \texttt{\$MaximumHistoryCount} -- command history size, 4096 by default.
\item \texttt{\$true} / \texttt{\$false} -- the boolean values.
\item \texttt{\$null} -- nothing. Use it to declare an empty variable, or to swallow unwanted output: \texttt{\$null = \$list.Add(x)}.
\item \texttt{\$LASTEXITCODE} -- exit code of the last external .exe (0 = success).
\end{itemize}

\textbf{Preference variables}

\begin{verbatim}
Get-Variable -Name *Preference
$ErrorActionPreference = 'SilentlyContinue'
$ErrorActionPreference = 'Continue'  # restore
\end{verbatim}
They control the output streams: Error, Warning, Verbose, Debug, Information.
See \texttt{Get-Help about\_Preference\_Variables}.

\section{5. Data Types}
PowerShell infers the type and lets a variable change type freely.
\begin{itemize}
\item \textbf{Bool}: \texttt{\$true} / \texttt{\$false}.
\item \textbf{Int32} (default, about $\pm 2.1$ billion), \textbf{Int64} for larger numbers.
\item \textbf{Double} (default for decimals), \textbf{Float} -- both can be imprecise.
\item \textbf{String}: must be quoted.
\end{itemize}
\begin{verbatim}
$num = 1
$num.GetType().name  # Int32
$num = 1.5
$num.GetType().name  # Double
[Int32]$num  # 2  (cast -> rounds)
\end{verbatim}

\textbf{Double vs single quotes}

\begin{verbatim}
$color = 'blue'
"I like $color"  # -> I like blue
'I like $color'  # -> I like $color
\end{verbatim}
Double quotes trigger \emph{variable interpolation} (expansion). Single quotes are literal.
With no quotes at all PowerShell reads the word as a command and errors out.

\section{6. Objects}
\textbf{Everything in PowerShell is an object.}
A class is the template, an object is one instance of it.
\textbf{Properties} are its data, \textbf{methods} are what it can do. Together they are its \emph{members}.
\begin{verbatim}
Select-Object -InputObject $color -Property *
Get-Member -InputObject $color  # props + methods
Get-Member -InputObject $color -Name Remove
$color.Length  # dot notation
$newColor = $color.Remove(1,1)  # methods take ()
\end{verbatim}
A method with a return type (e.g.\ \texttt{string Remove(...)}) produces a new object, so assign the result.
A \texttt{void} method modifies the existing object in place.

\section{7. Arrays}
\begin{verbatim}
$colorPicker = @('blue','white','yellow','black')
$colorPicker[0]  # blue  (indexes start at 0)
$colorPicker[-1]  # last element
$colorPicker[1..3]  # range operator ..
$colorPicker[3] = 'white'  # modify
$colorPicker += 'orange'  # add one
$colorPicker += @('pink','cyan')  # add several
\end{verbatim}
\textbf{Note}: arrays have a \emph{fixed size}. Every \texttt{+=} builds a brand new array, which hurts on large data. There is no simple way to remove an element.

\section{8. ArrayLists}
Dynamically sized, much faster than arrays for heavy adding/removing (roughly 100+ elements).
\begin{verbatim}
$cp = [System.Collections.ArrayList]@('blue','white')
$null = $cp.Add('gray')  # Add returns an index -> discard it
$cp.Remove('gray')  # removes by VALUE
\end{verbatim}

\section{9. Hashtables}
\textbf{Key--value} pairs, each key must be unique.
\begin{verbatim}
$users = @{
  abertram  = 'Adam Bertram'
  raquelcer = 'Raquel Cerillo'
}
$users['abertram']  # read
$users.abertram  # read as a property
$users.Keys ; $users.Values
$users.Add('natice','Natalie Ice')
$users['phrigo'] = 'Phil Rigo'  # add / modify
$users.ContainsKey('johnnyq')  # True/False
$users.Remove('natice')
\end{verbatim}

\section{10. Custom Objects}
\begin{verbatim}
$o = New-Object -TypeName PSCustomObject
$o = [PSCustomObject]@{ OSBuild='x'; OSVersion='y' }
$o.OSBuild
Get-Member -InputObject $o
\end{verbatim}
Build a hashtable whose keys are property names, then cast it to \texttt{[PSCustomObject]}. A \texttt{;} separates the pairs.

\section{11. The Pipeline}
The pipe operator \texttt{|} sends the \emph{output} of one command into the next as input. Unlike cmd.exe or bash, PowerShell pipes \textbf{objects}, not plain strings.
\begin{verbatim}
command1 | command2
Get-Service -Name 'wuauserv' | Start-Service
'wuauserv' | Get-Service | Start-Service
\end{verbatim}
An array sent down the pipeline is \textbf{unwrapped}: each item travels separately, so the next command runs once per item.
\begin{verbatim}
Get-Content -Path C:\Services.txt | Get-Service
Get-Content -Path C:\Services.txt | Get-Service | Start-Service
\end{verbatim}
No hard limit on chained commands, but past about five, rethink the approach. Not every command accepts pipeline input.

\section{12. Parameter Binding}
\textbf{Parameter binding} is how PowerShell decides which parameter receives the piped object. The command's author must explicitly enable it.
\begin{verbatim}
Get-Help -Name Get-Service -Full  # see PARAMETERS
\end{verbatim}
Look at the \texttt{Accept pipeline input?} field:
\begin{itemize}
\item \textbf{ByValue}: matched by object \emph{type}. Only one parameter per type can bind this way (e.g.\ a string binds to \texttt{-Name}).
\item \textbf{ByPropertyName}: matched by \emph{property name} on the incoming object (e.g.\ a \texttt{ComputerName} property binds to \texttt{-ComputerName}).
\end{itemize}
\begin{verbatim}
$o = [PSCustomObject]@{Name='wuauserv'; ComputerName='SRV1'}
$o | Get-Service
\end{verbatim}
Piping into a command with no pipeline support fails: \emph{The input object cannot be bound to any parameters}.

\section{13. Scripts}
A script is an external \texttt{.ps1} file holding a sequence of commands, run by typing its path.
\begin{verbatim}
C:\FolderPath\script.ps1  # full path
.\script.ps1  # current directory
\end{verbatim}
Write scripts in any editor; the PowerShell ISE (\texttt{powershell\_ise.exe}, now deprecated) or VS Code are the usual choices.
\begin{verbatim}
Write-Host 'Hello, I am in a script!'
\end{verbatim}

\textbf{Execution policy}

Scripts are blocked by default. Four settings:
\begin{itemize}
\item \textbf{Restricted} -- the default, no scripts at all.
\item \textbf{AllSigned} -- only scripts signed by a trusted party.
\item \textbf{RemoteSigned} -- your own scripts freely, downloaded ones only if signed.
\item \textbf{Unrestricted} -- any script.
\end{itemize}
\begin{verbatim}
Get-ExecutionPolicy
Set-ExecutionPolicy -ExecutionPolicy RemoteSigned
\end{verbatim}
Needs Administrator, but only once -- the value lives in the registry (or Group Policy across a domain). A \emph{script signature} is an encrypted comment block appended to the file.

\section{14. Comparison Operators}
Conditions are expressions evaluating to \texttt{True} / \texttt{False}.
\begin{itemize}
\item \texttt{-eq} equal, \texttt{-ne} not equal.
\item \texttt{-gt} greater, \texttt{-ge} greater or equal.
\item \texttt{-lt} less, \texttt{-le} less or equal.
\item \texttt{-contains} value is in an array.
\item \texttt{-not} inverts the result.
\end{itemize}
\begin{verbatim}
1 -eq 1  # True
Test-Connection -ComputerName srv1 -Quiet -Count 1
-not (Test-Connection -ComputerName srv1 -Quiet -Count 1)
\end{verbatim}
\texttt{-Quiet} reduces \texttt{Test-Connection} to a plain boolean, \texttt{-Count 1} to a single attempt.

\section{15. if / elseif / else}
\begin{verbatim}
if (condition) {
  # runs when the condition is True
}
\end{verbatim}
\texttt{\#} starts a comment. Add \texttt{else} for the false branch, and \texttt{elseif} for extra cases instead of nesting.
\begin{verbatim}
if (-not (Test-Connection -ComputerName $s -Quiet -Count 1)) {
  Write-Error -Message "$s is not responding!"
} elseif ($s -eq $problemServer) {
  Write-Error -Message "$s does not have the right file!"
} else {
  Get-Content -Path "\\$s\c$\App_configuration.txt"
}
\end{verbatim}
\texttt{elseif} branches are mutually exclusive: the first True branch wins and the rest are never tested.

\section{16. switch}
Cleaner than a long \texttt{elseif} chain when you are testing \emph{one value} against many possibilities.
\begin{verbatim}
switch ($currentServer) {
  $servers[0] {
    # do something
    break
  }
  $servers[1] {
    break
  }
  default {
    # no match was found
  }
}
\end{verbatim}
Unlike \texttt{elseif}, a \texttt{switch} keeps evaluating the remaining conditions after a match -- add \texttt{break} to make the branches mutually exclusive.

\section{17. Loops}
DRY -- Don't Repeat Yourself. Each pass of a loop is an \emph{iteration}. PowerShell has five loop types.

\textbf{foreach -- three flavors}

\begin{verbatim}
foreach ($server in $servers) {  # statement
  Get-Content -Path "\\$server\c$\f.txt"
}
$servers | ForEach-Object -Process {  # cmdlet, takes pipeline
  Get-Content -Path "\\$_\c$\f.txt"
}
$servers.foreach({  # method, PS v4+, fastest
  Get-Content -Path "\\$_\c$\f.txt"
})
\end{verbatim}
\texttt{\$\_} is the current pipeline object. The loop variable is only a \textbf{copy}, so no flavor of \texttt{foreach} can modify the original array.

\textbf{for -- a known number of times}

\begin{verbatim}
for ($i = 0; $i -lt 10; $i++) {
  $i
}
\end{verbatim}
Four parts: init, condition, step, body. Unlike \texttt{foreach}, it \emph{can} modify the array, and it can reach neighbouring elements.
\begin{verbatim}
for ($i = 0; $i -lt $servers.Length; $i++) {
  $servers[$i] = "new $($servers[$i])"
}
for ($i = 1; $i -lt $servers.Length; $i++) {
  Write-Host $servers[$i] "comes after" $servers[$i-1]
}
\end{verbatim}

\textbf{while, do/while, do/until}

\begin{verbatim}
while ($counter -lt 10) {  # 0 or more runs
  $counter
  $counter++
}
do {  # runs at least once
} while ($true)
do {
  $choice = Read-Host -Prompt 'Best language?'
} until ($choice -eq 'PowerShell')
\end{verbatim}
Use \texttt{while} when the iteration count is unknown. \texttt{do/while} and \texttt{do/until} are inverses of each other, and both test \emph{after} the body. \texttt{break} exits any loop -- essential inside \texttt{while (\$true)}.

\section{18. Errors \& Exceptions}
An \textbf{exception} is an event that disrupts normal flow (divide by zero, index out of bounds, missing file). Unhandled, it reaches the user as an \textbf{error}.
\begin{itemize}
\item \textbf{Terminating error} -- stops execution. Use when the code cannot continue at all (e.g.\ the whole folder is missing).
\item \textbf{Nonterminating error} -- prints a message and carries on. Use when one item fails but the rest are independent (e.g.\ one corrupt file out of many).
\end{itemize}
Whether an exception becomes terminating is a decision made by whoever wrote the cmdlet.
\begin{verbatim}
Get-Service bits,foo,lanmanserver
# errors on 'foo', still returns bits and lanmanserver
\end{verbatim}

\section{19. ErrorAction}
\texttt{-ErrorAction} is a \emph{common parameter}, available on every cmdlet. It says what to do about a \textbf{nonterminating} error:
\begin{itemize}
\item \texttt{Continue} -- print it and keep going (default).
\item \texttt{Ignore} -- no message, not recorded in \texttt{\$Error}.
\item \texttt{Inquire} -- print it and prompt the user.
\item \texttt{SilentlyContinue} -- no message, but still recorded in \texttt{\$Error}.
\item \texttt{Stop} -- print it and stop: turns it into a terminating error.
\end{itemize}
\begin{verbatim}
Get-ChildItem -Path $folderPath -ErrorAction Stop
$ErrorActionPreference = 'Stop'  # global default
\end{verbatim}
\texttt{-ErrorAction} overrides \texttt{\$ErrorActionPreference}. Setting the preference to \texttt{Stop} once removes the whole nonterminating category, so \texttt{try/catch} can see everything.

\section{20. try / catch / finally}
Only \textbf{terminating} errors can be caught -- hence the \texttt{-ErrorAction Stop} above.
\begin{verbatim}
try {
  # code to attempt
} catch {
  # runs only if a terminating error occurred
} finally {
  # always runs; optional
}
\end{verbatim}
Inside \texttt{catch}, the error object is in \texttt{\$\_}:
\begin{verbatim}
try {
  $files = Get-ChildItem -Path $folderPath -ErrorAction Stop
  $files.foreach({ (Get-Content $_)[0] })
} catch {
  $_.Exception.Message
}
\end{verbatim}
The error object also carries the exception type and a stack trace; \texttt{Message} is usually what you want. \texttt{finally} is the place for cleanup (closing database connections, removing remoting sessions).

\textbf{The \$Error variable}

\begin{verbatim}
$Error  # array of every error this session
$Error[0]  # most recent
$Error[1]  # second most recent
\end{verbatim}
Newest errors are added to the \emph{front}. See also \texttt{Get-Help about\_try\_catch\_finally}.

\section{21. Functions}
A \textbf{function} is a named block of code solving one task, defined once and called by name. A \textbf{cmdlet} does the same job but is compiled from C\#; functions are written in PowerShell.
\begin{verbatim}
function Install-Software {
  Write-Host 'I installed some software!'
}
Install-Software
Get-Command -CommandType Function
\end{verbatim}
Name them \texttt{Verb-Noun} (\texttt{Get-Verb} lists the approved verbs), noun singular. Redefining a function simply replaces it.

\section{22. Parameters}
Parameters live in a \texttt{param} block. \texttt{[CmdletBinding()]} or a \texttt{[Parameter()]} attribute marks an \emph{advanced} function.
\begin{verbatim}
function Install-Software {
  [CmdletBinding()]
  param(
    [Parameter(Mandatory)]
    [ValidateSet('1','2')]
    [string]$Version = 2
  )
  Write-Host "Installed version $Version."
}
Install-Software -Version 2
\end{verbatim}
\begin{itemize}
\item \texttt{[string]} casts the argument -- optional, but strongly recommended.
\item \texttt{Mandatory} makes PowerShell prompt when the value is missing; parameters are optional by default.
\item \texttt{= 2} sets a default value, overridden by anything passed in.
\item \texttt{ValidateSet} restricts the allowed values and fails early with a clear message. Many other validation attributes exist: \texttt{Get-Help about\_Functions\_Advanced\_Parameters}.
\end{itemize}

\section{23. Pipeline Input}
Add \texttt{ValueFromPipeline} (whole object, ByValue) or \texttt{ValueFromPipelineByPropertyName} (one property) to the parameter attribute.
\begin{verbatim}
function Install-Software {
  param(
    [Parameter(Mandatory)]
    [string]$Version,
    [Parameter(Mandatory, ValueFromPipeline)]
    [string]$ComputerName
  )
  process {
    Write-Host "Version $Version on $ComputerName."
  }
}
@('SRV1','SRV2') | Install-Software -Version 2
\end{verbatim}
Without a \texttt{process} block the body runs \emph{once}, for the last piped object only. \texttt{begin} and \texttt{end} blocks run before and after the whole pipeline: \texttt{Get-Help about\_Functions\_Advanced}.

\section{24. Modules}
A \textbf{module} is a group of related functions packaged for reuse -- at its simplest, a \texttt{.psm1} text file. Since v3 PowerShell auto-imports a module when one of its commands is called.
\begin{verbatim}
Get-Module  # imported in this session
Get-Module -ListAvailable  # installed on disk
Get-Command -Module Microsoft.PowerShell.Management
Import-Module -Name Software
Import-Module -Name Software -Force  # reload after edits
Remove-Module -Name Software  # unload, stays on disk
\end{verbatim}
Only \textbf{exported} commands are callable from outside; unexported ones are \emph{private} (helper) functions.

\textbf{Where modules live}

\begin{itemize}
\item \textbf{System}: \texttt{C:\textbackslash Windows\textbackslash System32\textbackslash WindowsPowerShell\textbackslash 1.0\textbackslash Modules} -- internal, don't add your own.
\item \textbf{All Users}: \texttt{C:\textbackslash Program Files\textbackslash WindowsPowerShell\textbackslash Modules}.
\item \textbf{Current User}: \texttt{...\textbackslash Documents\textbackslash WindowsPowerShell\textbackslash Modules}.
\end{itemize}
\begin{verbatim}
$env:PSModulePath  # ';' separated search paths
$env:PSModulePath + ';C:\MyModules'  # this session only
[Environment]::SetEnvironmentVariable("PSModulePath",
  $CurrentValue + ";C:\MyModules", "Machine")  # persistent
\end{verbatim}

\section{25. Module Anatomy \& Gallery}
A module is a \emph{folder} named after the module, holding a \texttt{.psm1} (the code) and an optional \texttt{.psd1} \textbf{manifest} -- a hashtable of metadata. Keep the noun identical across a module's functions; if the noun has to change, split the module.
\begin{verbatim}
mkdir 'C:\Program Files\WindowsPowerShell\Modules\Software'
Add-Content '...\Software\Software.psm1' -Value ''
New-ModuleManifest -Path '...\Software\Software.psd1' `
  -Author 'Adam Bertram' -RootModule Software.psm1 `
  -Description 'This module helps in deploying software.'
\end{verbatim}
At minimum set \texttt{RootModule}, \texttt{Author}, \texttt{Description}, \texttt{ModuleVersion}. \texttt{FunctionsToExport} controls exactly which commands are exported.

\textbf{PowerShell Gallery}

\texttt{PowerShellGet} is the built-in module for talking to \emph{powershellgallery.com}.
\begin{verbatim}
Get-Command -Module PowerShellGet
Find-Module -Name *VMware*  # search only
Find-Module -Name VMware.PowerCLI | Install-Module
Uninstall-Module -Name VMware.PowerCLI
\end{verbatim}
The Gallery is untrusted by default, so installing prompts; change it with \texttt{Set-PSRepository}. \texttt{Install-Module} defaults to the All Users path. \texttt{Remove-Module} only unloads -- \texttt{Uninstall-Module} deletes from disk, and only works on Gallery modules.

\vfill
\hrule
\smallskip
{\scriptsize Source: \emph{PowerShell for Sysadmins}, Ch.\ 1--7. Examples run on PowerShell v5.1.}
